\section{Reductions to canonical nonvanishing}
\label{red:section}

We first discharge the hypothesis of the projective theorem proved in
the appendices. We then reduce the remaining problem to canonical
nonvanishing on a smooth compact Kähler fourfold of irregularity zero.
Throughout, a fibration means a surjective holomorphic map with connected
fibers. Smooth source and target spaces do not mean that the map is
everywhere a submersion.

\subsection{The projective case}

\begin{lemma}[The logarithmic premise]\label{red:log-iitaka}
Assumption~\ref{hyp:campana} implies the following statement. Let
\(f:M\to S\) be a fibration between smooth connected complex projective
varieties, and let \(D_M,D_S\) be reduced effective simple normal crossing
divisors, possibly zero, such that
\begin{equation}\label{red:boundary-compatibility}
 \Supp(f^*D_S)\subseteq\Supp(D_M).
\end{equation}
For a very general smooth fiber \(F\), with \(D_F=D_M|_F\), one has
\begin{equation}\label{red:log-iitaka-inequality}
 \kappa(M,K_M+D_M)
 \geq \kappa(F,K_F+D_F)+\kappa(S,K_S+D_S).
\end{equation}
The usual convention for a \(-\infty\) summand applies.
\end{lemma}

\begin{proof}
It suffices to prove
\begin{equation}\label{red:invariant-base-bound}
 \kappa(f\mid D_M)\geq\kappa(S,K_S+D_S),
\end{equation}
because Assumption~\ref{hyp:campana} can then be applied to the original
pair \((M,D_M)\). We construct a neat model in the precise equivalence
class occurring in that assumption.

Flatten \(f\) by a projective modification of its base and resolve that
base, obtaining \(q:S'\to S\) with \(S'\) smooth projective. The main
strict transform before resolving the source can be taken
equidimensional over \(S'\). Indeed, start with the flat strict transform
provided by flattening and then base-change to the smooth resolved
base. Flatness persists; the unchanged smooth connected generic fiber
is irreducible, and flatness rules out vertical components. The main
space consequently has pure fibers of dimension \(\dim M-\dim S\).
Resolve it and its boundary to obtain a commuting diagram
\[
\begin{tikzcd}
 M' \arrow[r,"p"] \arrow[d,"f'"'] & M \arrow[d,"f"]\\
 S' \arrow[r,"q"'] & S .
\end{tikzcd}
\]
Here \(p\) is a projective birational morphism, \(M'\) is smooth
projective, and \(f'\) has connected fibers: its general fiber is
connected, so Stein factorization over the normal base \(S'\) has trivial
finite factor. Set \(D_{M'}\) equal to the strict transform of \(D_M\)
plus the full reduced \(p\)-exceptional divisor, and arrange that its
support is simple normal crossing.

This is an allowed orbifold modification of the source pair. It pushes
\(D_{M'}\) to \(D_M\). If a prime upstairs has positive pullback order
over a coefficient-one prime of \(D_M\), it is either that prime's
strict transform or a \(p\)-exceptional prime; in both cases its
coefficient upstairs is one. Thus the infinite-multiplicity condition
in the definition of an orbifold morphism is satisfied.

Moreover, \(p\) witnesses neatness. Let \(E\subset M'\) be a divisor
whose image in \(S'\) has codimension \(c\geq2\). In the equidimensional
main space its image has dimension at most
\[
 (\dim S-c)+(\dim M-\dim S)=\dim M-c.
\]
Its image in \(M\) therefore has codimension at least two, so \(E\) is
\(p\)-exceptional. Subsequent resolutions preserve this argument.

Let \(B_{f'}\) be the orbifold base divisor on \(S'\). Every component
of a pullback over a prime of
\((q^{-1}\Supp D_S)_{\mathrm{red}}\) belongs to \(D_{M'}\): this follows
from \eqref{red:boundary-compatibility} if it is not exceptional over
\(M\), and from the definition of \(D_{M'}\) otherwise. The multiplicity
assigned to each such source component is infinite. Consequently
\begin{equation}\label{red:orbifold-boundary-dominates}
 B_{f'}\geq(q^{-1}\Supp D_S)_{\mathrm{red}}.
\end{equation}
There is also an effective \(q\)-exceptional divisor \(E_S\) such that
\begin{equation}\label{red:log-base-comparison}
 K_{S'}+(q^{-1}\Supp D_S)_{\mathrm{red}}
   =q^*(K_S+D_S)+E_S
\end{equation}
for compatible canonical divisors. Above the boundary, this is the
nonnegativity of log discrepancies of the simple normal crossing pair
\((S,D_S)\); away from the boundary, it is the nonnegativity of the
ordinary discrepancies of the smooth space \(S\). Equivalently, a
blowup of a smooth center of codimension \(c\) contained in \(s\) boundary
components contributes \(c-s\) if the center lies in the boundary, and
\(c-1\) otherwise, both nonnegative.

The neat-model formula in Assumption~\ref{hyp:campana}, followed by
\eqref{red:orbifold-boundary-dominates} and
\eqref{red:log-base-comparison}, now gives
\[
 \kappa(f\mid D_M)
 =\kappa(S',K_{S'}+B_{f'})
 \geq\kappa(S,K_S+D_S).
\]
This proves \eqref{red:invariant-base-bound}. In particular, the infimum
in the definition of the invariant has been computed on an allowed neat
model, rather than replaced by the divisor of an arbitrary base model.
\end{proof}

\begin{corollary}[Projective abundance in the present setting]
\label{red:projective}
Under Assumption~\ref{hyp:campana}, a nef \(\Q\)-Cartier adjoint of a
normal projective log canonical pair over \(\C\), with effective rational
boundary, is semiample.
\end{corollary}

\begin{proof}
Lemma~\ref{red:log-iitaka} is exactly the logarithmic-Iitaka hypothesis
of Theorem~\ref{proj:conditional-log-abundance}. Apply that theorem,
whose full proof is included in the appendices.
\end{proof}

In particular, the Moishezon case of Theorem~\ref{thm:nonvanishing} is
settled. A compact Kähler Moishezon space with rational singularities is
projective by Namikawa's criterion \cite{NamikawaProjectivity}. Analytic klt
singularities are rational; this follows as well from relative vanishing
on a projective log resolution \cite{FujinoAnalyticMMP}. Thus a
Moishezon klt endpoint \((Y,\Delta)\) of
Assumption~\ref{hyp:mmp} is projective. Its rational analytic boundary
and adjoint are algebraic by Chow's theorem and GAGA, and analytic
nefness implies nonnegative degree on every curve. Corollary
\ref{red:projective} applies to its actual adjoint line.

\subsection{Lower-dimensional inputs and restriction to fibers}

We use smooth compact Kähler resolutions of spaces, pairs, and graphs.
The resolution and Kähler-modification theorems ensure that these can
be obtained by projective modifications and remain Kähler
\cite{HironakaResolution,VarouchasKahler}. A space in Fujiki class \(\mathcal C\)
has a smooth compact Kähler model. The quantities
\(\kappa(K_W)\), \(q(W)=h^0(W,\Omega_W^1)\), and \(a(W)\) for smooth
compact Kähler \(W\) are bimeromorphic invariants.

We recall precisely the lower-dimensional nonvanishing facts used
below. A smooth non-uniruled compact Kähler manifold of dimension at
most three has nonnegative canonical Kodaira dimension. In
nonprojective dimension three this is
\cite[Corollary~1.4]{HPMinimalModels}. More explicitly, take a terminal
minimal model by Höring--Peternell and apply canonical nonvanishing
\cite[Theorem~0.3]{DemaillyPeternellNonvanishing} to its nef canonical
class. Terminal discrepancy comparison pulls its plurisections back
to the smooth model.
In the projective case, the klt minimal model program and log abundance
in dimensions at most three imply nonvanishing for every effective
rational klt pair with pseudo-effective adjoint \cite{ThreefoldFlipsAbundance,KeelMatsukiMcKernan,KeelMatsukiMcKernanCorrection}.
Finally, Ou's Theorem~1.1 identifies non-uniruledness of a smooth compact
Kähler manifold with analytic pseudo-effectivity of its canonical class
\cite{OuUniruledness}. These inputs have no four-dimensional nonvanishing
conclusion.

We will repeatedly use the following elementary parameter observation.
Let \(f:Z\to S\) be a proper fibration between smooth compact complex
manifolds, and let a pseudo-effective rational adjoint on \(Z\) have a
closed positive representative with local potentials. Its restriction
to a smooth fiber is defined and positive for almost every parameter:
local plurisubharmonic potentials restrict unless they are identically
minus infinity, and Fubini excludes the latter event for almost every
parameter. This full-measure set meets the complement of any countable
union of proper analytic subsets. If the lower-dimensional results give
nonvanishing on those fibers, a fixed multiple has sections generically.
Indeed, on a connected smooth parameter open set the loci
\[
 \{s:h^0(Z_s,mL|_{Z_s})\geq1\},
\]
for divisible positive integers \(m\), are analytic jumping loci by
proper semicontinuity. If every one were proper, their union would have
measure zero. Some such locus is therefore the whole open set. Generic
base change gives a direct image of positive rank. The same reasoning
allows all required very-general smoothness and boundary conditions to
be imposed simultaneously.

\subsection{The uniruled and irregular cases}

\begin{lemma}[The uniruled reduction]\label{red:uniruled}
Assume Assumption~\ref{hyp:campana}.
Let \((Y,\Delta)\) be a non-Moishezon compact Kähler klt fourfold with
effective rational boundary and pseudo-effective \(\Q\)-Cartier adjoint
\(J=K_Y+\Delta\). If a smooth compact Kähler resolution of \(Y\) is
uniruled, then \(\kappa(Y,J)\geq0\).
\end{lemma}

\begin{proof}
Take the almost-holomorphic rational-chain quotient of a smooth
resolution \(W\) of \(Y\). Campana's rational quotient theorem in class
\(\mathcal C\) gives a quotient with base in that class and very general
fibers rationally connected after resolution
\cite[Theorem~2.6]{CampanaRationalQuotient}. Resolve the quotient and
the pair simultaneously to a fibration
\[
 f:W'\longrightarrow T,\qquad p:W'\longrightarrow Y,
\]
with \(W',T\) smooth compact Kähler. The very general fibers are
projective: rational-chain-connected manifolds in class \(\mathcal C\)
are Moishezon by Campana's algebraic connectedness criterion
\cite{CampanaAlgebraicConnectedness}, and a smooth Kähler Moishezon
manifold is projective. They are then rationally connected.

The quotient base is not uniruled; see also \cite[Lemma~8.10]{OuUniruledness}.
For completeness, suppose otherwise and choose a covering rational
curve through a very general point of \(T\). Pull the fibration back
to its normalization \(\PP^1\) and resolve the main component, obtaining
a smooth compact Kähler space \(Z\) over \(\PP^1\) with rationally
connected very general fiber. Such a fiber has no holomorphic one- or
two-forms. A global two-form on \(Z\) restricts to zero on those fibers;
the relative differential sequence over the smooth locus then places
it in the base-one-form times relative-one-form piece. Its restriction
there is again zero, so the two-form vanishes on a dense open and hence
everywhere. Thus \(H^0(Z,\Omega_Z^2)=0\). Rational approximation of a
Kähler class, followed by Kodaira's criterion, makes \(Z\) projective.
The Graber--Harris--Starr theorem supplies a section of
\(Z\to\PP^1\) \cite{GraberHarrisStarr}.

The chosen base point can be taken so that its original quotient fiber
is not contained in any of the countably many analytic exceptional
sets in the very-general quotient property. This follows from the
proper fiber-dimension theorem on the resolved quotient. Consequently
the rational curve is contained in none of the corresponding bad
parameter sets. Choose two distinct good parameters on it. In each
smooth projective rationally connected fiber, rational chains join a
very general point to the point of the section. The section joins
these two chains. Their projection to \(W\) is a rational chain joining
points of different very general rational-quotient fibers, where the
original almost-holomorphic quotient is defined. This contradicts the
quotient property. Hence \(T\) is not uniruled. Its dimension is
positive, since a point quotient would make \(W\) Moishezon, and is at
most three; the fiber dimension is likewise between one and three.

On the simultaneous log resolution write
\begin{equation}\label{red:klt-resolution}
 K_{W'}+\Gamma\sim_{\Q}p^*J+E,
 \qquad 0\leq\Gamma<1,\quad E\geq0,
\end{equation}
where \(\Gamma\) is simple normal crossing and \(E\) is
\(p\)-exceptional. Concretely, take the nonnegative part of the crepant
log pullback boundary; its negative part becomes \(E\).
The adjoint in \eqref{red:klt-resolution} is pseudo-effective. At a
very general smooth fiber \(F\) we can impose the klt simple normal
crossing restriction and also restrict its positive current, by the
parameter observation above. Projective nonvanishing in dimension at
most three gives \(\kappa(F,K_F+\Gamma|_F)\geq0\). The non-uniruled
base similarly satisfies \(\kappa(T,K_T)\geq0\).

For a smooth base, the invariant orbifold-base dimension is at least
its ordinary canonical Kodaira dimension: on every equivalent smooth
model the orbifold divisor is effective, and the smooth canonical
Kodaira dimension is birationally invariant. Assumption
\ref{hyp:campana} therefore gives
\(\kappa(W',K_{W'}+\Gamma)\geq0\). A resulting section pushes through
\eqref{red:klt-resolution} to a section of a positive multiple of
\(J\). The exceptional pole allowance disappears away from a
codimension-two subset of the normal target, and reflexive extension
fills that subset.
\end{proof}

\begin{lemma}[Irregular canonical nonvanishing]\label{red:irregular}
Assume Assumption~\ref{hyp:campana}. If \(W\) is a smooth non-uniruled
compact Kähler fourfold with \(q(W)>0\), then \(\kappa(W,K_W)\geq0\).
\end{lemma}

\begin{proof}
Resolve the Stein-factor base of the Albanese map and its main pullback.
This gives a fibration on smooth compact Kähler models, with a
positive-dimensional smooth base \(T\) mapping generically with full
rank to the Albanese torus. A suitable wedge of the invariant
one-forms on that torus pulls back to a nonzero section of \(K_T\).
The very general fiber has dimension at most three; if its dimension
is zero, it is a point. The pseudo-effective canonical class of the
total space restricts to the canonical class of almost every smooth
fiber. Lower-dimensional canonical nonvanishing and the parameter
observation supply \(\kappa(K_F)\geq0\) on very general fibers.
Assumption~\ref{hyp:campana}, with zero boundary, now gives canonical
nonvanishing on the resolved total space and hence on \(W\).
\end{proof}

\begin{proposition}[Canonical reduction]\label{red:reduction}
Under the three assumptions of the main theorem, it is enough for
Theorem~\ref{thm:nonvanishing} to prove canonical nonvanishing for every
smooth non-uniruled non-Moishezon compact Kähler fourfold \(W\) with
\(q(W)=0\).
\end{proposition}

\begin{proof}
Let \((Y,\Delta)\) be the nef klt endpoint in
Theorem~\ref{thm:nonvanishing}, and put \(J=K_Y+\Delta\). The
Moishezon case was settled by Corollary~\ref{red:projective}. Let
\(W\) be a smooth compact Kähler resolution in the remaining case.
If \(W\) is uniruled, Lemma~\ref{red:uniruled} applies. Otherwise
Ou's theorem makes \(K_W\) pseudo-effective. A section of a divisible
canonical power on \(W\) pushes to the corresponding reflexive
canonical power on \(Y\), and multiplication by the effective boundary
section, at a common multiple, gives a section of \(J\). Thus canonical
nonvanishing on \(W\) suffices. Lemma~\ref{red:irregular} handles
\(q(W)>0\), leaving exactly the stated case.

Once nonvanishing on this original endpoint is established,
Assumption~\ref{hyp:effective-abundance} makes its adjoint semiample.
Its other good-minimal-model and discrepancy properties have already
been supplied by Assumption~\ref{hyp:mmp}. None of the auxiliary
models constructed later needs to replace this endpoint.
\end{proof}
